\documentclass[10pt]{amsart}
\usepackage[a4paper,margin=23mm]{geometry}
\usepackage{amsmath,amssymb,amsthm,mathtools,hyperref,graphicx,xspace,letltxmacro,xpatch}
\usepackage[english]{babel}
\usepackage{xfrac}

\RequirePackage[mathscr]{eucal}
\usepackage{mathrsfs}
\let\mathcal \EuScript
\multlinegap0pt
\renewcommand{\theenumi}{\roman{enumi}}
\newcommand\mto{\mathchoice{\longmapsto}{\mapsto}{\mapsto}{\mapsto}}

\newcommand{\Psfrac}[2]{(\sfrac{#1}{#2})}
\newcommand{\psfrac}[2]{\sfrac{(#1)}{#2}}
\newcommand{\spfrac}[2]{\sfrac{#1}{(#2)}}

\newcommand{\scrL}{\mathscr{L}}
\newcommand{\scrO}{\mathscr{O}}
\newcommand{\scrQ}{\mathscr{Q}}

\newcommand{\cB}{\mathcal{B}}
\newcommand{\cC}{\mathcal{C}}
\newcommand{\cL}{\mathcal{L}}
\newcommand{\cQ}{\mathcal{Q}}
\newcommand{\cV}{\mathcal{V}}

\DeclareMathOperator{\dist}{dist}
\DeclareMathOperator{\Id}{Id}\let\sfId\Id
\DeclareMathOperator{\Span}{span}
\DeclareMathOperator{\Tr}{\mathsf{Tr}}

\newcommand{\intt}{\mathrm{int}}
\newcommand{\loc}{\mathrm{loc}}

\let\sep\mid

\providecommand{\enVert}[2][-1]{
\ensuremath{\mathinner{
\ifthenelse{\equal{#1}{-1}}{ % if
\!\left\lVert#2\right\rVert}{}
\ifthenelse{\equal{#1}{0}}{ % if
\lVert#2\rVert}{}
\ifthenelse{\equal{#1}{1}}{ % if
\!\bigl\lVert#2\bigr\rVert}{}
\ifthenelse{\equal{#1}{2}}{ % if
\!\Bigl\lVert#2\Bigr\rVert}{}
\ifthenelse{\equal{#1}{3}}{ % if
\!\biggl\lVert#2\biggr\rVert}{}
\ifthenelse{\equal{#1}{4}}{ % if
\!\Biggl\lVert#2\Biggr\rVert}{}
}}
}
\let\norm=\enVert

\usepackage{csquotes,dsfont,braket}

\theoremstyle{plain}
\newtheorem{theorem}{Theorem}[section]
\newtheorem{lemma}[theorem]{Lemma}
\newtheorem{proposition}[theorem]{Proposition}
\newtheorem{corollary}[theorem]{Corollary}
\theoremstyle{definition}
\newtheorem{notation}[theorem]{Notation}
\newtheorem{definition}[theorem]{Definition}
\newtheorem{remark}[theorem]{Remark}

\newcommand{\RR}{\mathbb{R}}
\newcommand{\CC}{\mathbb{C}}
\newcommand{\NN}{\mathbb{N}}
\newcommand{\dx}{\mathrm{\,d}}
\newcommand{\bn}{{\boldsymbol n}}
\newcommand{\bx}{{\boldsymbol x}}
\let\pa\partial
\let\eps\varepsilon
\DeclareMathOperator{\IM}{Im}\let\Im\IM
\DeclareMathOperator{\RE}{Re}\let\Re\RE
\newcommand{\Mat}{\mathsf{Mat}}
\newcommand{\jac}{\mathsf{Jac}}
\DeclareMathOperator{\supp}{\mathsf{supp}}
\DeclareMathOperator{\dive}{div}
\DeclareMathOperator{\Dom}{\mathsf{Dom}}


\usepackage{xfrac,xcolor,graphicx}\graphicspath{{./figures/}}\providecommand{\skpt}{\par\smallskip\noindent}\providecommand{\pointir}{.\ }\providecommand{\eg}{e.g.}\providecommand{\ie}{i.e.}\providecommand{\resp}{resp.}\providecommand{\cf}{cf.}\providecommand{\etc}{etc.}\let\Subsection\subsection\let\Subsubsection\subsubsection

\usepackage{ifthen}
